\documentclass[11pt,a4paper]{article}

\usepackage[utf8]{inputenc}
\usepackage[T1]{fontenc}
\usepackage[spanish,es-nodecimaldot]{babel}
\usepackage{lmodern}
\usepackage[margin=2.5cm]{geometry}
\usepackage{amsmath,amssymb}
\usepackage{booktabs,tabularx,array}
\usepackage{enumitem}
\usepackage{microtype}
\usepackage{xurl}
\usepackage[hidelinks]{hyperref}
\usepackage{fancyhdr}

\setlength{\headheight}{14pt}
\pagestyle{fancy}
\fancyhf{}
\fancyhead[L]{Propuestas de ampliación del yugo FFB}
\fancyhead[R]{Proyecto Final}
\fancyfoot[C]{\thepage}
\setlength{\parindent}{0pt}
\setlength{\parskip}{0.5em}
\setlist[itemize]{itemsep=2pt,topsep=4pt}
\setlist[enumerate]{itemsep=2pt,topsep=4pt}
\newcolumntype{Y}{>{\raggedright\arraybackslash}X}

\title{Propuestas de ampliación y validación experimental\\
de un yugo con realimentación de fuerza}
\author{Proyecto Final de Ingeniería Electrónica\\
Universidad Tecnológica Nacional --- Facultad Regional Córdoba}
\date{Plan de estudios 2023}

\begin{document}
\maketitle

\begin{abstract}
Se presentan cinco líneas de trabajo complementarias para un yugo de dos ejes
con realimentación de fuerza: medición del esfuerzo aplicado al usuario,
instrumentación de corriente, estimación de velocidad, caracterización
electromecánica y supervisión de fallas con integración electrónica.
Se conserva la numeración de las opciones originalmente propuestas: 1, 2, 3,
4 y 7. Cada opción identifica su relación con la carrera, el desarrollo
previsto y los resultados experimentales necesarios para justificarla.
\end{abstract}

\tableofcontents
\newpage

\section*{Alcance y criterios generales}
\addcontentsline{toc}{section}{Alcance y criterios generales}

El sistema considerado reproduce efectos de fuerza constante y sinusoidal en
los ejes de alabeo (\textit{roll}) y cabeceo (\textit{pitch}). La posición puede
obtenerse mediante potenciómetros o encoders incrementales A/B. Se contempla
la incorporación de un lazo PI de corriente y una contribución de
amortiguación calculada a partir de la velocidad medida.

Las propuestas se vinculan con asignaturas incluidas en el Plan 2023 de
Ingeniería Electrónica de la UTN FRC~\cite{plan2023}. Esta vinculación es una
propuesta de integración académica; el alcance definitivo deberá acordarse
con la dirección del Proyecto Final.

Este documento describe actividades por desarrollar y validar. No acredita
que el control de corriente, las protecciones o los ensayos aquí planteados
se encuentren implementados. Las magnitudes de fuerza, corriente, velocidad
y temperatura admisibles deberán definirse a partir del hardware definitivo.
Los criterios de aceptación se establecerán antes de realizar los ensayos,
evitando seleccionarlos a posteriori para ajustarlos a los resultados.

\section{Medición de la fuerza real aplicada al piloto}
\label{sec:fuerza}

\subsection{Objetivo y vinculación curricular}
Determinar la relación entre el comando solicitado, la corriente del motor
y el esfuerzo efectivamente aplicado al yugo. Esta actividad permite integrar
Medidas Electrónicas I y II, Electrónica Aplicada, Física I y Probabilidad y
Estadística.

El PWM representa una acción de mando eléctrica y no constituye por sí mismo
una medición de fuerza. Incluso con corriente regulada, la transmisión,
fricción y geometría pueden introducir diferencias entre el torque del motor
y el esfuerzo que percibe el usuario.

\subsection{Instrumentación propuesta}
Se propone utilizar una celda de carga con acondicionamiento y adquisición
adecuados. Puede integrarse a un banco de ensayo desmontable, sin incorporarse
obligatoriamente al lazo de control final.

En \textit{pitch}, la celda deberá medir la componente de fuerza longitudinal
en la dirección de desplazamiento. En \textit{roll}, se aplicará una fuerza
a una distancia conocida del eje. El torque se calculará mediante
\begin{equation}
    \tau_r = F\,\ell\,\sin\alpha,
\end{equation}
donde $F$ es la fuerza medida, $\ell$ es la distancia al eje y $\alpha$ el
ángulo entre el brazo y la fuerza. Si el montaje garantiza perpendicularidad,
la expresión se reduce a $\tau_r=F\ell$.

La calibración inicial podrá representarse mediante
\begin{equation}
    \widehat F = a\,n+b,
\end{equation}
donde $n$ es la lectura digital y $a$, $b$ se obtienen con cargas de referencia.
Se verificará experimentalmente si el modelo lineal resulta suficiente.

\subsection{Metodología experimental}
\begin{enumerate}
    \item Calibrar cero y sensibilidad del sistema de medición, verificando
    ambos sentidos de carga y distintos puntos del rango.
    \item Aplicar referencias escalonadas de corriente dentro del rango
    autorizado para el banco. Registrar corriente, fuerza, posición y tiempo.
    \item Repetir el ensayo en diferentes posiciones del recorrido y con
    secuencias ascendentes y descendentes para evaluar histéresis.
    \item Realizar repeticiones en condiciones comparables de temperatura
    y montaje para estimar dispersión.
    \item Incorporar ensayos dinámicos solamente si la celda, el
    acondicionamiento y la adquisición poseen ancho de banda suficiente.
\end{enumerate}

Para una referencia física de fuerza $F^*[k]$ y $N$ muestras, podrá utilizarse
\begin{equation}
    \operatorname{RMSE}_F =
    \sqrt{\frac{1}{N}\sum_{k=1}^{N}
    \left(\widehat F[k]-F^*[k]\right)^2}.
\end{equation}
La referencia deberá expresarse en newtons mediante una conversión
documentada; los valores abstractos del protocolo FFB no equivalen
directamente a una fuerza física.

Para el caso perpendicular y despreciando correlaciones, una aproximación
de la incertidumbre estándar combinada del torque es
\begin{equation}
    u_c(\tau_r)=\sqrt{\ell^2u^2(F)+F^2u^2(\ell)}.
\end{equation}
Si la alineación no es suficientemente precisa, se incluirá también la
contribución angular y cualquier correlación relevante.

\subsection{Resultados esperados}
Se obtendrán curvas corriente--fuerza y corriente--torque, estimaciones de
histéresis y repetibilidad, y un presupuesto de incertidumbre. El resultado
permitirá evaluar con unidades físicas el desempeño del yugo y contrastar
el modelo desarrollado en la opción~\ref{sec:modelo}.

\section{Diseño y calibración de la medición de corriente}

\subsection{Objetivo y vinculación curricular}
Diseñar una cadena de medición que permita realimentar el PI de corriente
con resolución, ancho de banda y error conocidos. Se relaciona con
Electrónica Aplicada I y II, Teoría de los Circuitos, Medidas Electrónicas
y Electrónica de Potencia. La medición de corriente es un elemento central
de las aplicaciones de control de torque de motores DC~\cite{tiBDC}.

\subsection{Alternativas de implementación}
\begin{tabularx}{\textwidth}{@{}p{0.25\textwidth}YY@{}}
\toprule
Alternativa & Aspectos a evaluar & Limitaciones a verificar \\
\midrule
Shunt y amplificador & Sensibilidad, ganancia, costo y ancho de banda &
Disipación, tensión de modo común y ubicación del shunt \\
Sensor Hall & Medición bidireccional y aislamiento, si el modelo lo ofrece &
Offset, deriva térmica, ruido y sensibilidad \\
Salida de medición del driver & Integración y cantidad de componentes &
Escala, estados de validez, ancho de banda y error especificado \\
\bottomrule
\end{tabularx}

Para un shunt $R_s$ y un amplificador de ganancia $G$, el modelo ideal es
\begin{equation}
    V_o=V_0+G R_s i,
    \qquad
    \widehat i=\frac{V_o-V_0}{G R_s},
\end{equation}
donde $V_0$ es el nivel de salida correspondiente a corriente nula.
La disipación del shunt se estima mediante
\begin{equation}
    P_s=I_{\mathrm{rms}}^2R_s.
\end{equation}
Para un ADC ideal de $b$ bits y referencia $V_{\mathrm{ref}}$, el incremento
aproximado de corriente por cuenta es
\begin{equation}
    \Delta i \approx \frac{V_{\mathrm{ref}}}{2^bG R_s}.
\end{equation}
Esta resolución nominal no contempla ruido, no linealidad ni incertidumbre
de la referencia.

\subsection{Desarrollo propuesto}
\begin{enumerate}
    \item Definir corriente máxima, polaridad, resolución requerida y
    ancho de banda compatible con el control.
    \item Seleccionar el sensor y dimensionar ganancia, offset y margen
    frente a saturación del amplificador y del ADC.
    \item Diseñar el filtro analógico y analizar su atenuación y retardo.
    Para un RC de primer orden, $f_c=1/(2\pi RC)$; su suficiencia deberá
    evaluarse frente al espectro de ruido y a la frecuencia de muestreo.
    \item Calibrar la cadena completa con una referencia adecuada en
    varios puntos y, cuando corresponda, en ambos sentidos de corriente.
    \item Observar la señal con PWM activo y verificar qué se mide durante
    conducción, recirculación y frenado. La corriente tomada de la fuente
    no representa necesariamente la corriente del devanado en todo instante.
    \item Evaluar adquisición sincronizada con PWM cuando la topología y
    el microcontrolador lo permitan. El instante de muestreo debe justificarse
    según los intervalos válidos de medición~\cite{microchipSync}.
\end{enumerate}

\subsection{Resultados esperados}
El desarrollo entregará esquemático, cálculos, curva de calibración,
resolución efectiva, error y deriva observados. Deberá quedar documentada
la función de transferencia de la medición y su influencia sobre el PI.

\section{Estimación de velocidad y procesamiento de señales}

\subsection{Objetivo y vinculación curricular}
Obtener una estimación de velocidad adecuada para la amortiguación local,
reduciendo ruido sin introducir un retardo excesivo. Se aplican Análisis de
Señales y Sistemas, Teoría de los Circuitos II y Técnicas Digitales.

\subsection{Fundamentos}
Tanto los potenciómetros como los encoders permiten estimar el sentido de
movimiento a partir de la variación temporal de posición. Una aproximación
por diferencias finitas es
\begin{equation}
    \widehat v[k]=\frac{x[k]-x[k-1]}{t[k]-t[k-1]}.
\end{equation}
La diferenciación hace más visible el ruido y la cuantización de posición.
Por ello, se propone comparar esta estimación con una versión filtrada y
con una estimación realizada sobre una ventana de muestras.

Un filtro de primer orden puede implementarse mediante
\begin{align}
    v_f[k]&=\alpha_k v_f[k-1]+(1-\alpha_k)\widehat v[k],\\
    \alpha_k&=\frac{\tau_f}{\tau_f+\Delta t_k},
\end{align}
donde $\tau_f$ es la constante de tiempo y $\Delta t_k=t[k]-t[k-1]$.
Esta forma corresponde a una discretización por Euler hacia atrás del
filtro continuo $H(s)=1/(1+\tau_f s)$. Permite explicitar el compromiso entre
atenuación de ruido y retardo.

Para una coordenada lineal, la fuerza de amortiguación ideal es
\begin{equation}
    F_d=-B_v v_f,
    \qquad [B_v]=\mathrm{N\,s/m}.
\end{equation}
Para una coordenada angular, se utilizará
\begin{equation}
    \tau_d=-B_\omega\omega_f,
    \qquad [B_\omega]=\mathrm{N\,m\,s/rad}.
\end{equation}
Si el firmware utiliza posición y comando normalizados, sus coeficientes
de amortiguación no tienen directamente estas unidades físicas. La
conversión deberá documentarse antes de comparar ambos modelos.

\subsection{Metodología experimental}
\begin{itemize}
    \item Registrar señales en reposo y durante movimientos lentos y rápidos,
    manteniendo la misma calibración física para cada comparación.
    \item Comparar potenciómetros y encoders considerando resolución,
    repetibilidad, ruido y respuesta ante cambios de sentido.
    \item Medir desviación estándar de velocidad en reposo, retardo y
    respuesta dinámica de cada estimador. Para cuantificar error de
    velocidad se necesitará una referencia independiente o una trayectoria
    de banco conocida; no basta con comparar dos estimadores entre sí.
    \item Registrar intervalos de adquisición reales y evaluar su variación.
    La frecuencia nominal del lazo no garantiza muestreo uniforme.
    \item Comparar el retorno del yugo con y sin amortiguación, utilizando
    condiciones iniciales y referencias de fuerza reproducibles.
\end{itemize}

\subsection{Resultados esperados}
Se justificarán la frecuencia de adquisición, el estimador de velocidad y
los parámetros del filtro mediante curvas temporales y métricas. También
se cuantificará el efecto de la amortiguación sobre el sobrepaso y las
oscilaciones, sin asumir que una mejora subjetiva demuestra estabilidad.

\section{Caracterización del motor y de la transmisión}
\label{sec:modelo}

\subsection{Objetivo y vinculación curricular}
Construir y validar un modelo electromecánico para explicar la conversión
entre corriente, torque y fuerza en el yugo. Esta opción integra Física I
y II, Máquinas e Instalaciones Eléctricas, Análisis de Señales y Sistemas
y Sistemas de Control.

\subsection{Modelo inicial}
Suponiendo un motor DC de imanes permanentes con escobillas, se plantea
\begin{align}
    v_a &= R_a i_a + L_a\frac{di_a}{dt}+K_e\omega_m,\\
    \tau_m &= K_t i_a,\\
    J_{\mathrm{eq}}\frac{d\omega_m}{dt}
        &= \tau_m-B_{\mathrm{eq}}\omega_m-\tau_f-\tau_{L,m}.
\end{align}
Aquí $v_a$ es la tensión en bornes del motor, $R_a$ y $L_a$ son los parámetros
eléctricos, $K_e$ y $K_t$ sus constantes electromecánicas,
$J_{\mathrm{eq}}$ y $B_{\mathrm{eq}}$ representan parámetros referidos al eje
motor, $\tau_f$ agrupa fricción no viscosa y $\tau_{L,m}$ es el torque de carga
referido a ese eje. La validez del modelo se comprobará para el motor utilizado.

Definiendo una reducción $n=\omega_m/\omega_p$, para operación motriz
cuasiestática puede aproximarse
\begin{equation}
    \tau_p \approx \eta n\tau_m,
    \qquad
    F_p\approx\frac{\tau_p}{r_p},
\end{equation}
donde $\eta$ es la eficiencia y $r_p$ el radio efectivo de la polea de pitch.
La operación inversa, cuando el usuario mueve el motor, puede requerir
otro modelo de pérdidas. No se supondrá la misma eficiencia en todos los
sentidos, velocidades y cargas sin contrastación experimental.

Para una polea dentada con $z$ dientes y paso $p$, el avance ideal de correa
por vuelta es $zp$. Si un encoder de $P$ pulsos por vuelta se encuentra
solidario a dicha polea y se decodifica en cuadratura $\times4$,
\begin{equation}
    x=\frac{q}{4P}\,zp,
\end{equation}
donde $q$ es el desplazamiento acumulado en cuentas. Si existe otra relación
entre encoder y polea, deberá incorporarse explícitamente.

\subsection{Procedimiento de identificación}
\begin{enumerate}
    \item Registrar datos de placa, geometría, relación de transmisión y
    condiciones de los instrumentos utilizados.
    \item Estimar resistencia e inductancia mediante ensayos eléctricos
    limitados en corriente y duración, registrando la temperatura.
    \item Estimar las constantes electromecánicas a partir de mediciones
    de velocidad, tensión, corriente y esfuerzo, según la instrumentación
    disponible y las condiciones del ensayo.
    \item Caracterizar fricción, holgura y elasticidad mediante movimientos
    en ambos sentidos y diferentes niveles de carga.
    \item Ajustar los parámetros y validar el modelo con registros distintos
    de los empleados para identificarlo.
\end{enumerate}

\subsection{Resultados esperados}
Se obtendrán parámetros con unidades e incertidumbres o intervalos estimados,
curvas medidas y simuladas, y un dominio de validez del modelo. Estos
resultados permitirán fundamentar la elección de motor, transmisión y
geometría, y explicar las diferencias entre roll y pitch.

\setcounter{section}{6}
\section{Supervisión de fallas e integración electrónica}

\subsection{Objetivo y vinculación curricular}
Definir el comportamiento del sistema ante condiciones anómalas y desarrollar
una interconexión electrónica reproducible. Se vincula con Seguridad,
Higiene y Medio Ambiente, Técnicas Digitales y Tecnología Electrónica.

La supervisión por software se complementará con las medidas de hardware
que correspondan al riesgo del equipo. Un microcontrolador bloqueado no
puede garantizar la ejecución de su propia rutina de parada.

\subsection{Máquina de estados propuesta}
\begin{tabularx}{\textwidth}{@{}p{0.23\textwidth}Y@{}}
\toprule
Estado & Comportamiento previsto \\
\midrule
Inicialización & Salidas de potencia deshabilitadas y comprobación inicial
de sensores y configuración. \\
Calibración & Obtención de referencias de posición y de cero de corriente;
las maniobras deberán definirse para el montaje. \\
Preparado & Mediciones disponibles y potencia aún no habilitada. \\
Habilitado & Ejecución del control con límites y supervisión activos. \\
Falla & Aplicación de la respuesta definida y registro de la causa.
La reactivación requiere eliminar la causa y un rearme explícito. \\
\bottomrule
\end{tabularx}

\subsection{Condiciones a supervisar}
\begin{itemize}
    \item Sobrecorriente: umbral instantáneo y, si se justifica, criterio
    temporal asociado al calentamiento.
    \item Sobretemperatura del driver o motor, con histéresis para evitar
    conmutaciones repetidas alrededor del umbral.
    \item Valores de posición o corriente fuera de rango y comprobaciones
    de plausibilidad. Un encoder inmóvil no demuestra por sí solo una falla,
    y una entrada analógica desconectada puede entregar valores aparentemente
    válidos; la cobertura deberá evaluarse con el circuito real.
    \item Recorrido excesivo y actuación de finales de carrera, si se incorporan.
    \item Parada de emergencia, reinicio inesperado y vigilancia de ejecución
    mediante \textit{watchdog}.
    \item Pérdida de comunicación, diferenciándola de la ausencia legítima
    de actualizaciones de un efecto persistente.
\end{itemize}

En particular, un efecto CONSTANT puede permanecer activo sin paquetes
continuos. Por lo tanto, un temporizador reiniciado solamente por recepción
de efectos no constituye una detección suficiente de desconexión.

La respuesta de falla debe especificar si corresponde deshabilitar la etapa,
reducir controladamente la referencia o aplicar otra acción. PWM cero,
frenado y desconexión no son estados eléctricamente equivalentes para todos
los drivers; deberá verificarse la tabla de funcionamiento del dispositivo
y la energía almacenada en el sistema.

\subsection{Placa de adquisición e interconexión}
Se propone una placa que integre conexiones de sensores, acondicionamiento,
distribución de alimentación y señales de habilitación. El diseño deberá
documentar retornos de corriente, desacoplo, protección de entradas,
conectores, puntos de prueba y separación funcional entre señales sensibles
y trayectos de potencia. La estrategia de masas se justificará según los
caminos reales de corriente y no mediante una separación arbitraria.

\subsection{Verificación y entregables}
Se elaborará una matriz que relacione falla, mecanismo de detección, respuesta,
tiempo máximo admisible y condición de rearme. Se ensayarán fallas simuladas
en condiciones controladas, inicialmente con la potencia deshabilitada o
limitada. Se registrarán tiempos de respuesta, falsas detecciones y fallas
que el sistema no logra detectar.

Los entregables incluirán máquina de estados, matriz de fallas, esquemáticos,
diseño de PCB cuando corresponda y resultados de ensayo. Estas actividades
no equivalen a una certificación de seguridad del producto.

\section*{Secuencia sugerida de desarrollo}
\addcontentsline{toc}{section}{Secuencia sugerida de desarrollo}

Antes de los ensayos con potencia, se definirán las medidas básicas de
limitación y parada de la opción 7. Luego se recomienda calibrar la medición
de corriente (opción 2) y el banco de fuerza (opción 1), caracterizar el
conjunto electromecánico (opción 4) y evaluar los estimadores y filtros
(opción 3). La integración final permitirá comparar el mando PWM de
referencia, el PI de corriente y el PI con amortiguación local bajo
condiciones reproducibles.

Para cada comparación se documentarán versión del firmware, parámetros,
calibraciones, instrumentos, condiciones iniciales y archivos de datos.
La evidencia central será la concordancia entre requisitos, cálculos,
simulación y mediciones, incluyendo las limitaciones observadas.

\begin{thebibliography}{9}
\bibitem{plan2023}
Universidad Tecnológica Nacional, Facultad Regional Córdoba.
\textit{Plan de estudio de la carrera de Ingeniería Electrónica, Plan 2023}.
\url{https://www.institucional.frc.utn.edu.ar/electronica/pub/file/Plan%202023/Plan%20de%20estudio%20-%20Lista%20de%20asignaturas%202023.pdf}.

\bibitem{tiBDC}
Texas Instruments.
\textit{Brushed DC motor drivers: current sensing and control}.
\url{https://www.ti.com/product-category/motor-drivers/brushed-dc/overview.html}.

\bibitem{microchipSync}
Microchip Technology.
\textit{Hardware Timing and Synchronization}.
Ejemplo de sincronización PWM--ADC en control de motores; su aplicación
al yugo requiere adaptación al hardware seleccionado.
\url{https://onlinedocs.microchip.com/oxy/GUID-69278EBA-9F00-4CEE-8103-2998236856BD-en-US-2/GUID-16339F99-D30F-48D1-8854-21B6454F6CEC.html}.
\end{thebibliography}

\end{document}
